\section{Simulating a Poisson r.v.}

Given a sequence $U_1,U_2,\ldots$ of i.i.d. $U[0,1]$
random variables, the following method simulates a
$\operatorname{Poi}(\lambda)$ r.v.

$\textbf{Simulation Method 7.}$

\begin{enumerate}[label=\arabic*.]
\item Sample the $U_k$'s one at a time until the first time
$m$ when $M_m =\prod_{k=1}^m U_k$ satisfies
$M_m \le e^{-\lambda}$.

\item Output $m-1$.
\end{enumerate}

\begin{xexercise}
Explain why the output $Z$ of this algorithm has
distribution $\operatorname{Poi}(\lambda)$, using the
connection between the Poisson distribution and the
cumulative sums of i.i.d. exponential random variables
(see Section \ref{sec:exp-dist}).
\end{xexercise}

\noindent Note that the number of samples this method requires
is equal to one plus the random variable $Z$ being
simulated. In particular, the mean number of samples
required is $\mathbf{E} Z = \lambda+1$.
